\documentclass[10.5pt,a4paper]{article}
\def\yo{C}
% Macros y tabla de orden compartidas por los tres guiones. \yo = A, B o C.
\usepackage[utf8]{inputenc}
\usepackage[T1]{fontenc}
\usepackage{lmodern}
\usepackage[spanish,es-nodecimaldot]{babel}
\usepackage[a4paper,margin=2cm]{geometry}
\usepackage{amsmath,amssymb}
\usepackage[table]{xcolor}
\usepackage{booktabs,longtable,array}
\usepackage{fancyhdr}
\usepackage{enumitem}
\usepackage[hidelinks]{hyperref}
\setlist{itemsep=1pt,topsep=2pt,parsep=0pt}
\setlength{\parindent}{0pt}
\setlength{\parskip}{4pt}
\definecolor{azulfila}{HTML}{DCE9F8}
\definecolor{hipo}{HTML}{FFF3D6}
\definecolor{nuevo}{HTML}{DFF3E3}
\newcommand{\dt}{\ensuremath{\Delta t}}
\newcommand{\dts}{\ensuremath{\Delta t^{*}}}
\newcommand{\vv}{\ensuremath{\langle v^4\rangle/\langle v^2\rangle^2}}
\newcommand{\kbt}{\ensuremath{k_BT}}
\newcommand{\tn}[2]{\ensuremath{#1\times10^{#2}}}
\newcommand{\pr}[1]{\par\textbf{P:} #1\par\nobreak}
\newcommand{\rp}[1]{\textbf{R:} #1\par}
\newcommand{\hipotesis}[1]{\par\colorbox{hipo}{\parbox{\dimexpr\linewidth-2\fboxsep}{\textbf{Hipótesis nuestra (no medida):} #1}}\par}
\newcommand{\nuevo}[1]{\par\colorbox{nuevo}{\parbox{\dimexpr\linewidth-2\fboxsep}{\textbf{Nuevo en la presentación:} #1}}\par}
\newcommand{\debil}[1]{\par\textbf{Punto débil:} #1\par}
\usepackage{xspace}
\newcommand{\pc}{\char37\xspace}

\pagestyle{fancy}
\fancyhf{}
\renewcommand{\headrulewidth}{0pt}
\newcommand{\mi}[1]{\cellcolor{azulfila}\textbf{#1}}
\newcommand{\Fmi}[5]{\mi{#1} & \mi{#2} & \mi{#3} & \mi{#4} & \mi{#5} \\}
\newcommand{\Fno}[5]{#1 & #2 & #3 & #4 & #5 \\}
\newcommand{\F}[5]{\ifnum\pdfstrcmp{#3}{\yo}=0 \expandafter\Fmi\else\expandafter\Fno\fi{#1}{#2}{#3}{#4}{#5}}
\newcommand{\interv}[8]{%
  \par\vspace{6pt}\noindent{\large\bfseries Diapositiva #1 \textemdash{} #2}\hfill{\small reloj #3 \textbullet{} #4 s}\par\nobreak
  {\small\textit{Entrás después de:} #5}\par\nobreak
  #6\par
  \ifx\relax#7\relax\else\colorbox{hipo}{\parbox{\dimexpr\linewidth-2\fboxsep}{\small\textbf{Números a mano:} #7}}\par\fi
  {\small\textit{#8}}\par}
\newcommand{\nota}[1]{\colorbox{nuevo}{\parbox{\dimexpr\linewidth-2\fboxsep}{\small #1}}\par}
\newcommand{\cabecera}[3]{%
  {\LARGE\bfseries Guion Orador #1}\par
  {\large TP4 \textemdash{} Dinámica molecular regida por el paso temporal \textemdash{} Presentación de 13 minutos}\par\vspace{2pt}
  \textbf{Tu tiempo hablando:} #2 en #3 intervenciones. \textbf{Tiempo total planeado de la presentación:} 12:30 (límite 13:00, quedan 30 s de margen para pausas y pases).

  Los textos están pensados a unas 140--150 palabras por minuto. Si te quedás corto o largo, respetá los segundos de la tabla: el reloj de cada diapositiva es acumulado. Los números de diapositiva son los del pie «n / 33» de la presentación vigente.}
\newcommand{\tablaorden}{%
\section{Orden completo (tus filas en azul)}
{\small
\begin{longtable}{@{}rp{8.6cm}crr@{}}
\toprule Diap. & Contenido & Voz & Seg. & Reloj \\ \midrule \endhead
\F{1}{Título}{A}{10}{0:00}
\F{3}{Sistema real}{A}{20}{0:10}
\F{4}{Modelo matemático}{B}{25}{0:30}
\F{6}{Integración: Verlet original}{C}{25}{0:55}
\F{7}{Arquitectura del motor}{A}{25}{1:20}
\F{9}{Sistema simulado}{B}{25}{1:45}
\F{10}{Observables (1/2)}{C}{30}{2:10}
\F{11}{Observables (2/2)}{A}{30}{2:40}
\F{12}{Corridas realizadas}{B}{20}{3:10}
\F{14}{Oscilador: ECM contra $\Delta t$}{C}{30}{3:30}
\F{15}{Energía total en función del tiempo}{A}{20}{4:00}
\F{16}{Error de energía contra $\Delta t$}{B}{40}{4:20}
\F{17}{Tiempo de ejecución en función de $N$}{C}{40}{5:00}
\F{18}{Tiempo de ejecución en función de la densidad (nueva)}{B}{20}{5:40}
\F{19}{Variación de $x_0$ (animaciones)}{A}{20}{6:00}
\F{20}{Fracción de usadas en función del tiempo}{B}{25}{6:20}
\F{21}{Tiempo al 90\,\pc contra $x_0$ (parte 1)}{C}{30}{6:45}
\F{21}{Tiempo al 90\,\pc contra $x_0$ (parte 2)}{A}{20}{7:15}
\F{22}{Billar sin obstáculos (animación)}{A}{15}{7:35}
\F{23}{Estacionario: $k_BT$ ajustado en el tiempo (nueva)}{B}{25}{7:50}
\F{24}{Distribución de rapideces: evolución}{C}{25}{8:15}
\F{25}{Búsqueda del mejor ajuste}{A}{25}{8:40}
\F{26}{Distribución estacionaria y ajuste}{B}{30}{9:05}
\F{27}{Tiempos de conversión contra la densidad}{C}{30}{9:35}
\F{28}{$\langle t_{90}\rangle$ y $\langle t_{100}\rangle$ contra la densidad (nueva)}{A}{20}{10:05}
\F{29}{Fracción de usadas a $t_f$ contra la densidad}{A}{20}{10:25}
\F{30}{Mapa de $t_{90}$ en $(x_0,N)$ --- parte 1}{B}{25}{10:45}
\F{30}{Mapa de $t_{90}$ en $(x_0,N)$ --- parte 2}{C}{20}{11:10}
\F{32}{Conclusiones (1 a 3)}{A}{20}{11:30}
\F{32}{Conclusiones (4 a 6)}{B}{20}{11:50}
\F{32}{Conclusiones (7 y 8) y cierre}{C}{20}{12:10}
\bottomrule
\end{longtable}}

\textbf{Reparto.} Las diapositivas de resultados (la parte más importante) las hablamos los tres, intercaladas: A, B y C se turnan. En las diapositivas de tiempo al 90\,\pc contra $x_0$ (C y A) y del mapa de calor (B y C) hablan dos voces seguidas en la misma diapositiva; el pase se hace sin cambiar de lámina. Las diapositivas 14 a 18 no se detienen en «por qué»: las razones del método (Verlet, regla de $\Delta t^*$) ya no están escritas, se dicen de viva voz.

\textbf{Cambios respecto del guion anterior.} Las tres diapositivas nuevas son la 18, la 23 y la 28 (marcadas «nueva»); la antigua «Relajación hacia Maxwell-Boltzmann» (cociente) ya no tiene diapositiva propia: ahora su idea se dice en la 11 (definición) y en la 23 ($k_BT(t)$). Desde la 18 la numeración se corrió.}
\newcommand{\antes}{%
\section{Antes de subir}
Leé tu guion en voz alta con cronómetro: cada intervención debe entrar en sus segundos. Los pases son frases cortas (``sigue B'', ``te paso a C'') sin cambiar de tema.

Para las preguntas, mirá la \textbf{Guía de estudio TP4}: cada uno defiende sus diapositivas, pero cualquiera puede ser interrogado sobre cualquier parte.}

\cfoot{Guion Orador C \textbullet{} TP4 Grupo 5 \textbullet{} \thepage}
\begin{document}
\cabecera{C}{4:10}{9}
\tablaorden

\section{Tu guion}

\interv{6}{Integración: Verlet original}{0:55}{25}{B (Modelo matemático)}{Integramos con Verlet original, como pide el enunciado. La posición nueva sale de dos veces la actual, menos la anterior, más $\Delta t$ cuadrado sobre $m$ por la fuerza. La velocidad se obtiene por diferencia centrada, con un paso de atraso. Para arrancar estimamos la posición en menos $\Delta t$ con Taylor. El $\Delta t$ es fijo y guardamos el estado cada $n$ pasos.}{$x(t+\Delta t)=2x(t)-x(t-\Delta t)+\Delta t^2F/m$; $v(t)=[x(t+\Delta t)-x(t-\Delta t)]/2\Delta t$}{Pasás a A (Arquitectura del motor).}

\interv{10}{Observables (1/2)}{2:10}{30}{B (Sistema simulado)}{Todos los observables se calculan después, en Python; el motor sólo escribe estados. La energía suma la cinética más el potencial de cada contacto, incluida la pared. El error de energía $\varepsilon$ es el desvío relativo medio respecto de la energía inicial. La fracción de usadas es $N_u$ sobre $N$, y $t_{90}$ es el primer instante en que llega a 0,9. Promediamos sobre realizaciones, con el desvío estándar muestral como barra de error, y si alguna realización no llega al 90 por ciento lo informamos: nunca promediamos sólo las exitosas.}{barras $=\sigma$ muestral ($n-1$)}{Pasás a A (Observables 2/2).}

\interv{14}{Oscilador: ECM contra $\Delta t$}{3:30}{30}{B (Corridas realizadas)}{Antes del billar validamos los integradores con el oscilador amortiguado, que tiene solución exacta. Es el error cuadrático medio contra $\Delta t$, en log-log. Euler predictor-corrector tiene pendiente 2: orden 1. Beeman, Verlet original y Velocity Verlet tienen pendiente 4, es decir orden 2, y casi el mismo error. Elegimos Verlet original porque evalúa la fuerza una sola vez por paso y es el que pide el enunciado para el billar.}{pendientes 1,95 / 4,00; ECM($10^{-3}$ s): Euler PC $3\times10^{-4}$ m$^2$, Verlet $3{,}8\times10^{-10}$ m$^2$}{\textbf{Ojo:} «Método elegido: Verlet original» y sus razones ya no están escritas en la diapositiva; decilas. Pasás a A (Energía total en función del tiempo).}

\interv{17}{Tiempo de ejecución en función de $N$}{5:00}{40}{B (Error de energía contra $\Delta t$)}{Comparamos el tiempo de ejecución contra el del TP3, dirigido por eventos, en la misma computadora y con $t_f$ igual a 30 segundos. En log-log, TP4 tiene pendiente 1,01: crece lineal con $N$ gracias a la grilla de celdas, porque la cantidad de pasos es fija. TP3 tiene pendiente 3,62: en cada evento actualiza todas las partículas y los eventos crecen con la densidad. Con pocas partículas gana TP3; el cruce está en $N$ cercano a 310. En $N$ igual 650, TP4 tarda unos 8 segundos y TP3 unos 140.}{pendientes 1,01 (TP4) y 3,62 (TP3); cruce $N\approx310$; $N=650$: 7,9 s vs 145 s}{Pasás a B (Tiempo de ejecución en función de la densidad).}

\interv{21}{Tiempo al 90\,\pc contra $x_0$ (parte 1)}{6:45}{30}{B (Fracción de usadas en función del tiempo)}{Este es el resultado principal: $t_{90}$ contra $x_0$, para $N$ igual 100 y $N$ igual 20, con 10 realizaciones. Con $N$ igual 100 hay una zona favorable entre 0,15 y 0,35 metros, con $t_{90}$ de unos 15 segundos, y los extremos son peores: unos 24 segundos con los obstáculos tocándose y unos 31 pegados a la pared.}{$N=100$: $(14{,}6\pm1{,}8)$ s en 0,20 m; $(24\pm2)$ s en $x_0=r$; $(31\pm5)$ s en $R-r$}{Pasás a A (parte 2, misma lámina).}

\interv{24}{Distribución de rapideces: evolución}{8:15}{25}{B (Estacionario: $k_BT$ ajustado en el tiempo)}{Esta es $f$ de $v$ en el tiempo. En $t$ igual 0 es una delta en $v_0$, la flecha. En menos de medio segundo el pico se ensancha hasta la forma de equilibrio: aparecen partículas lentas y rápidas, algunas más rápidas que $v_0$, porque los choques intercambian energía y sólo se conserva el total.}{}{Pasás a A (Búsqueda del mejor ajuste).}

\interv{27}{Tiempos de conversión contra la densidad}{9:35}{30}{B (Distribución estacionaria y ajuste)}{Pasamos al efecto de la densidad con las corridas de $x_0$ igual a $r$ y $t_f$ igual 30 segundos, $N$ de 10 a 150. $t_{90}$ es mínimo para $N$ entre 20 y 40: $17\pm4$ segundos en $N$ igual 30, indistinguible de sus vecinos; es una zona, no un punto. Para $N$ más alto algunas realizaciones no llegan al 90 por ciento: graficamos una cota inferior, con símbolo hueco. Y $t_{100}$ queda censurado en todos los $N$: ninguno tiene las 10 realizaciones antes de $t_f$.}{$N=30$: $(17\pm4)$ s; $N=20$: $(19\pm4)$ s; $N=40$: $(20\pm3)$ s; $N=10$: 7 de 10 llegan}{Pasás a A (diapositiva 28, $\langle t_{90}\rangle$ y $\langle t_{100}\rangle$).}

\interv{30}{Mapa de $t_{90}$ en $(x_0,N)$ --- parte 2}{11:10}{20}{B (misma diapositiva)}{Para $N$ de 50 o más, el mínimo cae siempre en $x_0$ entre 0,20 y 0,25 metros: la posición favorable es casi independiente de la densidad. Lo que cambia es el valor: el $t_{90}$ mínimo sube de unos 11 segundos con $N$ igual 20 a unos 45 con $N$ igual 400, y los extremos se penalizan más cuanto más denso: con $N$ igual 400, $x_0$ igual a $r$ tarda el doble que 0,20.}{mínimo 0,20--0,25 m; $t_{90}$ mín: 11, 14, 15, 17, 20, 24, 29, 45 s ($N=20,\dots,400$); $N=400$: $(89\pm5)$ s en $x_0=r$ vs $(45\pm3)$ s en 0,20}{Pasás a A (Conclusiones 1 a 3).}

\interv{32}{Conclusiones (7 y 8) y cierre}{12:10}{20}{B (Conclusiones 4 a 6)}{Séptimo: $t_{90}$ es mínimo para $N$ entre 20 y 40 y crece con la densidad, junto con $t_{100}$: no hay una densidad óptima distinguible; la fracción de usadas baja al aumentar la densidad. Octavo: en el mapa, la posición óptima de los obstáculos es casi independiente de $N$. Muchas gracias; quedamos atentos a sus preguntas.}{}{Fin de la presentación: abrir preguntas.}

\antes
\end{document}
